\chapter{全同粒子与交换统计}
\label{ch:identical}

\section{全同粒子与交换对称性}
\chapref{ch:em}给出了一个 Landau 能级中的单电子态数。要从态数推得该能级可容纳的电子数，
还需确定两个电子能否占据同一个态。电子的电荷、质量和内禀自旋相同。
即使探测器能区分它们所占的模式，也不能给两个电子赋予可观测的粒子编号。

考虑两个互不重叠的空间模式 $a,b$，并设探测器能够分辨粒子到达的模式。
若入射的是电子和质子，$\ket{a}_{\electron}\otimes\ket{b}_{\mathrm p}$和
$\ket{b}_{\electron}\otimes\ket{a}_{\mathrm p}$可以由“电子在哪一端”区分。
换成两个电子，且自旋相同，实验仍可读出$a,b$各有一个电子，记录中却没有粒子编号。
为构造两粒子态，先给张量积的两个因子临时编号 $1,2$，按这个次序写
$\ket{a,b}=\ket a_1\otimes\ket b_2$。两个位置各填一个单粒子态。
\begin{definition}{全同粒子与交换对称性}{qm-exchange}
    全同粒子(identical particles)具有相同的内禀性质，交换其编号保持所有物理预测不变。
    交换算符(exchange operator) $\op{P}_{12}$ 把乘积基矢 $\ket{a,b}$ 变为 $\ket{b,a}$，并按线性作用于叠加态。
    对称态 $\ket{\Psi_+}$ 和反对称态 $\ket{\Psi_-}$ 分别满足
    \begin{equation}
        \op P_{12}\ket{\Psi_+}=+\ket{\Psi_+},\qquad
        \op P_{12}\ket{\Psi_-}=-\ket{\Psi_-}.
    \end{equation}
\end{definition}
探测器$a,b$仍各在原处。连续交换两次回到原态，故若一次交换只乘一个相位，
这个相位的平方为 1。两种符号分别对应哪类粒子，由自旋统计定理确定。

\begin{theorem}{自旋与统计的联系}{qm-spin-statistics}
    在三维空间的相对论量子场论中，在局域性、正能量与态空间正定等通常条件下，
    全同的整数自旋粒子服从 Bose 统计，称为 Bose 子；全同的半整数自旋粒子服从 Fermi 统计，称为 Fermi 子。
    相应的多粒子态对任意一对粒子编号的交换分别对称、反对称。
\end{theorem}
本章把自旋统计定理(spin--statistics theorem)的结论作为非相对论模型的输入。
定理的场论条件见\href{https://web2.ph.utexas.edu/~vadim/Classes/2024f-qft/spinstat.pdf}{得克萨斯大学的自旋统计讲义}。
二维准粒子的交换还可能产生其他相位，称为任意子(anyon)统计的情形需另行处理。

\section{全同粒子态的归一化与 Pauli 不相容原理}
交换算符的本征值决定了两个粒子能否共同占据某个单粒子态。
从两个不同的单粒子态出发，对交换前后的乘积态取对称或反对称组合，就可以考察两态重合时的极限。

设$\ket{a},\ket{b}$各自归一，重叠为$s=\braket{a}{b}$。
分子中的空间轨道、自旋态或光学脉冲模式都可以作为单粒子标签。
交换前后的两个乘积态一般有重叠，
\begin{equation}
    (\bra{a}\otimes\bra{b})(\ket{b}\otimes\ket{a})
    =\braket{a}{b}\braket{b}{a}=\abs{s}^2.
\end{equation}
因此未归一化对称或反对称组合的范数平方为
$1+1\pm\abs{s}^2\pm\abs{s}^2=2(1\pm\abs{s}^2)$。

相应的归一化态为
\begin{equation}
    \ket{a,b}_\pm=
    \frac{\ket{a,b}\pm\ket{b,a}}{\sqrt{2(1\pm\abs{s}^2)}}.
\end{equation}
只在$s=0$时分母才是$\sqrt2$。
\begin{corollary}[Pauli 不相容原理]
    两个全同 Fermi 子不能占据同一个完整的单粒子态。此时 $\abs{s}=1$，交换前后的两项在反对称组合中相消，
    得到的零向量不能归一化。两个电子仍可占据同一空间轨道，只要它们的自旋态正交。
\end{corollary}

\begin{definition}{Slater行列式}{qm-slater}
    将 $N$ 个正交归一单粒子态 $\ket{a_1},\ldots,\ket{a_N}$ 各占据一次，反对称化并归一化所得的态称为 Slater 行列式，
    \begin{equation}
        \ket{a_1,\ldots,a_N}_{\antisym}=\frac1{\sqrt{N!}}
        \sum_{\pi\in S_N}\sgn(\pi)\ket{a_{\pi(1)}}\otimes\cdots\otimes\ket{a_{\pi(N)}}.
    \end{equation}
    $S_N$包含$N$个编号的全部排列，$\sgn(\pi)$在偶排列时为$+1$、奇排列时为$-1$。
\end{definition}
两两正交的$N!$项有相同模长，因而归一化系数为$1/\sqrt{N!}$。
各个 Slater 行列式可作多体基矢；相互作用体系的本征态通常是多个行列式的叠加。
对于两个电子，反对称自旋单态要乘对称空间态，对称自旋三重态则要乘反对称空间态。

\section{Fock空间与占据数表示}
Slater 行列式用粒子编号写出所有排列。实际计数实验直接给出每个可分辨模式中的粒子数。
以这些占据数标记量子态，也便于描述光子发射等改变粒子总数的过程。

\begin{definition}{占据数与Fock空间}{qm-fock-space}
    选定一组正交归一的单粒子模式并固定次序，$\ket{n_1,n_2,\ldots}$表示各模式的粒子数分别确定为$n_1,n_2,\ldots$的占据数态。
    各模式占据数之和为 $N$ 的这些态张成 $N$ 粒子扇区。不同总粒子数的扇区彼此正交。
    汇集所有扇区的基便得到 Fock 空间的一组基，这种空间组合称为正交直和。
    用 $\ket{\Psi_N}$ 记总粒子数为 $N$ 的分量，归一化态的扇区分解满足
    \begin{equation}
        \ket{\Psi}=\sum_{N=0}^\infty\ket{\Psi_N},\qquad
        \sum_{N=0}^\infty\norm{\Psi_N}^2=1.
    \end{equation}
\end{definition}
例如只保留模式 $a,b$ 时，真空扇区有一个基矢 $\ket{0,0}$，单粒子扇区有两个基矢 $\ket{1,0},\ket{0,1}$，
合起来是 $1+2=3$ 维。再纳入两粒子扇区，便可描述两个模式各有一个粒子的 $\ket{1,1}$。
\tabref{tab:two-mode-occupations}列出只保留两个正交单粒子模式时的两粒子基。
表的第一列按模式 $a,b$ 的顺序记粒子数，后两列按临时编号 $1,2$ 的顺序记单粒子态。
例如 $\ket{1,1}$ 表示两个模式各有一个粒子，将同一占据情况写成带粒子编号的态时，须按交换对称性组合 $\ket{a,b}$ 和 $\ket{b,a}$。
这里的模式标签包含自旋等全部自由度。

\begin{table}[htbp]
    \centering
    \caption{两个全同粒子在两个正交模式中的占据数基。两个模式的占据数之和固定为2。}
    \label{tab:two-mode-occupations}
    \begin{tabular}{ccc}
        \toprule
        占据数态 $\ket{n_a,n_b}$ & 用粒子编号写的 Bose 态 & 用粒子编号写的 Fermi 态 \\
        \midrule
        $\ket{2,0}$ & $\ket{a,a}$ & 无允许态 \\
        $\ket{1,1}$ & $(\ket{a,b}+\ket{b,a})/\sqrt2$ & $(\ket{a,b}-\ket{b,a})/\sqrt2$ \\
        $\ket{0,2}$ & $\ket{b,b}$ & 无允许态 \\
        \bottomrule
    \end{tabular}
\end{table}

例如在模式$a,b$的单粒子能量分别为$\epsilon_a,\epsilon_b$时，
无相互作用的$\ket{1,1}$的能量就是$\epsilon_a+\epsilon_b$。
若往空模式逐个加入粒子，计数改变便对应态空间中的一个操作。
真空$\ket{\vac}$表示所有模式都空，产生算符$\op{a}_r\adjoint$在第$r$个模式加一个粒子，
湮灭算符$\op{a}_r$从这个模式取走一个粒子。给定一组正交单粒子模式$\ket{r}$，Bose 子的算符满足
\begin{equation}
    [\op{a}_r,\op{a}_s\adjoint]=\delta_{rs}\one,\quad [\op{a}_r,\op{a}_s]=0,\qquad
    \ket{n_1,n_2,\ldots}=
    \prod_r\frac{(\op{a}_r\adjoint)^{n_r}}{\sqrt{n_r!}}\ket{\vac}.
\end{equation}
Fermi 子的产生、湮灭算符满足反对易关系
\begin{equation}
    \{\op{c}_r,\op{c}_s\adjoint\}=\delta_{rs}\one,\qquad
    \{\op{c}_r,\op{c}_s\}=0.
\end{equation}
第二式取 $s=r$ 再取伴随，得到 $(\op{c}_r\adjoint)^2=0$，故 $n_r$ 只能取 $0$ 或 $1$。
写多Fermi子态时固定模式顺序；对第$r$个模式湮灭会带前面已占据数决定的符号。

\begin{definition}{占据数态上的产生与湮灭}{qm-fock-action}
    固定模式次序$1,2,\ldots$，用
    $\ket{\bm n}=\ket{n_1,n_2,\ldots}$表示占据数态(occupation-number state)。
    对Bose子，
    \begin{subequations}
        \begin{alignat}{1}
            \op{a}_r\ket{\bm n} &=\sqrt{n_r}\ket{\ldots,n_r-1,\ldots},\\
            \op{a}_r\adjoint\ket{\bm n} &=\sqrt{n_r+1}\ket{\ldots,n_r+1,\ldots}.
        \end{alignat}
    \end{subequations}
    对 Fermi 子，记排在第 $r$ 个模式之前的所有模式的占据数之和为 $N_{<r}$，则
    \begin{subequations}
        \begin{alignat}{1}
            \op{c}_r\ket{\bm n} &=(-1)^{N_{<r}}n_r\ket{\ldots,0_r,\ldots},\\
            \op{c}_r\adjoint\ket{\bm n} &=(-1)^{N_{<r}}(1-n_r)\ket{\ldots,1_r,\ldots}.
        \end{alignat}
    \end{subequations}
    真空 $\ket{\vac}$ 满足所有 $n_r=0$。总粒子数算符 $\op N$ 在占据数态上的本征值是各模式粒子数之和 $N$，即
    \[
        \op N=\sum_r \op{a}_r\adjoint \op{a}_r\quad\text{（Bose 子）},\qquad
        \op N=\sum_r \op{c}_r\adjoint \op{c}_r\quad\text{（Fermi 子）}.
    \]
\end{definition}

把 Fermi 子湮灭算符 $\op{c}_r$ 移到第 $r$ 个产生算符的位置，每经过一个已占据模式就出现一次负号。
例如，从前两个模式各有一个粒子的态中移去第二个粒子，必须交换一次算符，
\[
    \op{c}_2\op{c}_1\adjoint\op{c}_2\adjoint\ket{\vac}
    =-\op{c}_1\adjoint\op{c}_2\op{c}_2\adjoint\ket{\vac}
    =-\op{c}_1\adjoint\ket{\vac}.
\]
对 Bose 子，$\op{a}_1\adjoint\ket{1,1}=\sqrt2\ket{2,1}$，产生算符使两粒子态进入三粒子扇区。
给定$N$且演化守恒粒子数时，整个计算都在同一扇区内完成。

设单体算符(one-body operator)为 $\op{h}$，对称的两体相互作用(two-body interaction)为 $\op{v}$。多体 Hamilton 算符在占据数表示中为
\begin{equation}
    \op{H}=\sum_{rs}h_{rs}\op{a}_r\adjoint \op{a}_s
    +\frac12\sum_{rstu}V_{rs;tu}\op{a}_r\adjoint \op{a}_s\adjoint \op{a}_u \op{a}_t,
    \quad V_{rs;tu}=\mel{r,s}{\op{v}}{t,u}.
\end{equation}
Bose与Fermi粒子都采用这个表达式，Fermi情形把$\op{a}$换成$\op{c}$，按反对易关系运算。
$h_{rs}=\mel{r}{\op{h}}{s}$描述一个粒子从$s$模式到$r$模式的矩阵元。
两体项中的 $\op{a}_u\op{a}_t$ 在模式 $t,u$ 中各湮灭一个粒子，
$\op{a}_r\adjoint\op{a}_s\adjoint$ 则在模式 $s,r$ 中各产生一个粒子。
这里 $V_{rs;tu}$ 是在乘积基中计算的原始矩阵元，例如 $\ket{t,u}=\ket t_1\otimes\ket u_2$。求和同时包含两种粒子排列，故有系数 $1/2$。
在固定 $N$ 扇区，这对应于按粒子编号写出的 Hamilton 算符
\[
    \op{H}=\sum_{i=1}^N \op{h}(i)+\sum_{1\leq i<j\leq N}\op{v}(i,j).
\]

若把交换项吸收到反对称化矩阵元中，对全部指标求和时系数相应变为$1/4$。
无相互作用时，将每个模式的能量 $\epsilon_r$ 乘以占据数 $n_r$ 后相加，即得总能量。Pauli 原理决定 Fermi 子如何填充各个模式。

\section{Hong--Ou--Mandel符合凹陷}
\label{sec:hom}
上节的 $\ket{1,1}$ 可以通过在分束器的两个输入端各制备一个光子来实现。
为检验交换对称性如何影响输出，Hong--Ou--Mandel 实验在两个输出端分别探测光子，
改变两个输入的相对延迟$\tau$，统计一次试验中两端各记录一个光子的符合计数(coincidence counting)。
远离时间重叠时，两光子到达不同输出的符合率给出基线；
在频率、偏振和空间模式也匹配时，零延迟附近出现符合凹陷\cite{hom}。
总单计数不必同时消失，因此符合减少表示光子对输出端口的分配改变。

“两光子都透射”和“两光子都反射”通向同一符合结果。先按分束器的透反射相位叠加这两条过程的振幅，再取模平方，便得到符合概率。
以 $a,b$ 标记输入端口，$c,d$ 标记输出端口。取无损平衡分束器的相位约定，使相应的产生算符满足
\begin{equation}
    \op{a}\adjoint\longmapsto\frac{\op{c}\adjoint+\op{d}\adjoint}{\sqrt2},\qquad
    \op{b}\adjoint\longmapsto\frac{\op{c}\adjoint-\op{d}\adjoint}{\sqrt2}.
\end{equation}
两个输入端各一个、其他自由度完全不可分辨的光子满足
\begin{equation}
    \begin{alignedat}{1}
        \ket{1,1}_{ab}=\op{a}\adjoint \op{b}\adjoint\ket{\vac} &\longmapsto\frac{(\op{c}\adjoint)^2-(\op{d}\adjoint)^2}{2}\ket{\vac}\\
        &=\frac{\ket{2,0}_{cd}-\ket{0,2}_{cd}}{\sqrt2}.
    \end{alignedat}
\end{equation}
两输出各一个的振幅相消，符合概率为零。
如果两个光子只有部分重叠，符合计数会恢复多少？
设它们的归一化内部模式为 $\ket{u},\ket{v}$，忽略损耗，探测器只分辨输出端口。
先保留内部标签，对 $\op{a}\adjoint(u)\op{b}\adjoint(v)\ket{\vac}$ 作同一分束变换，
恰好每个输出一个光子的部分是
\begin{equation}
    \ket{\Psi_{1,1}}
    =\frac12\bigp{\ket{v}_c\otimes\ket{u}_d-\ket{u}_c\otimes\ket{v}_d}.
\end{equation}
下标$c,d$指输出端口。这个未归一化分量的范数平方就是符合概率：
\begin{equation}
    \begin{alignedat}{1}
        P_{1,1}=\braket{\Psi_{1,1}}{\Psi_{1,1}} &=\frac14\bigb{2-\braket{v}{u}\braket{u}{v}-\braket{u}{v}\braket{v}{u}}\\
        &=\frac12(1-\abs{\braket{u}{v}}^2).
    \end{alignedat}
\end{equation}
内部模式相同时，两个到达不同端口的振幅相消，符合概率为零；
内部模式正交时，交叉项为零，符合概率回到 $1/2$。
相对延迟通过单光子内部态的时间演化改变重叠：
若$\op{H}_\mathrm{ph}\ket{\omega}=\hbar\omega\ket{\omega}$，
则$\ket{v_\tau}=\exp(-\i \op{H}_\mathrm{ph}\tau/\hbar)\ket{v}$。
对相同频谱包络的纯单光子输入，记
$f(\omega)=\braket{\omega}{u}=\braket{\omega}{v}$，
把频谱归一化到总概率为 1，得到
\begin{equation}
    \braket{u}{v_\tau}
    =\int \abs{f(\omega)}^2\e^{-\i\omega\tau}\d\omega.
\end{equation}
这里$f$是单光子态在频率基中的分量，$\omega$是角频率。

\begin{example}{Gauss 频谱的符合曲线}{qm-hom-gaussian}
    设谱强度在载频 $\omega_0$ 附近为 Gauss 分布，
    角频率标准差为$\sigma_\omega\ll\omega_0$。
    将窄带积分延伸到整个实轴，有
    \begin{equation}
        \abs{f(\omega)}^2
        =\frac{1}{\sqrt{2\pi}\sigma_\omega}
            \exp\Bigp{-\frac{(\omega-\omega_0)^2}{2\sigma_\omega^2}}.
    \end{equation}
    Gauss 积分给出
    $\braket{u}{v_\tau}=\e^{-\i\omega_0\tau}
        \e^{-\sigma_\omega^2\tau^2/2}$，于是
    \begin{equation}
        P_{1,1}(\tau)
        =\frac12\bigp{1-\e^{-\sigma_\omega^2\tau^2}}.
    \end{equation}
    符合凹陷的时间尺度是$1/\sigma_\omega$；
    因而带宽越窄，两态的时间重叠就能维持到越大的延迟（图\ref{fig:hom-dip}）。
    偏振不重合或频谱混合会使零延迟处的重叠小于1，因而抬高最低符合率。
\end{example}

\begin{figure}[htbp]
    \centering\includegraphics[width=.60\linewidth]{figures/hom_dip.pdf}
    \caption{理想平衡分束器的符合概率 $P_{1,1}$。每条曲线的两个入射光子频谱相同、偏振平行；横轴用同一参考带宽 $\sigma_{\mathrm{ref}}$ 标度延迟，可在相同符合概率处比较两个凹陷的宽度。}
    \label{fig:hom-dip}
\end{figure}
这项分析适用于输入内部态可分解为两个纯单光子态的情形；
若光子对在频率上纠缠，应对联合频率态计算交换重叠。

\checkpoint{两电子的自旋单态反对称，
    对空间态施加了什么限制？}{
    设总态可分为$\ket{\Phi}\otimes\ket{s}$，
    其中 $\ket{s}=(\ket{\up,\dn}-\ket{\dn,\up})/\sqrt2$，两个箭头依次给出编号 $1,2$ 的自旋投影方向。
    交换两个电子时，空间标签和自旋标签必须同时交换。
    因为$\op{P}_{12}^{\spinlabel}\ket{s}=-\ket{s}$，
    要使总态满足$\op{P}_{12}\ket{\Psi}=-\ket{\Psi}$，
    就必须有$\op{P}_{12}^{\spatial}\ket{\Phi}=+\ket{\Phi}$。

    对两个正交空间轨道 $\ket{a},\ket{b}$，自旋单态要求取对称空间因子 $\ket{\Phi_+}$，自旋三重态要求取反对称空间因子 $\ket{\Phi_-}$。两者可合写为
    \begin{equation}
        \ket{\Phi_\pm}
        =\frac{\ket{a,b}\pm\ket{b,a}}{\sqrt2}.
    \end{equation}
    当两个电子占据同一空间轨道时，反对称组合的两项相消，
    所以只有自旋单态组合可以保留。

    这个限制还会改变相互作用能。
    对交换对称的两体势$\op{v}$，令
    $D=\mel{a,b}{\op{v}}{a,b}$、
    $K=\mel{a,b}{\op{v}}{b,a}$。
自伴性与交换对称性使 $K$ 为实数。将对称和反对称态分别代入势能期望值，得到
    \begin{equation}
        \mel{\Phi_\pm}{\op{v}}{\Phi_\pm}=D\pm K.
    \end{equation}
    同一个相互作用在对称和反对称空间态上的期望值相差 $2K$；这个交换能差来自两种粒子排列的振幅干涉。
    若空间与自旋本来就纠缠，应直接检查完整联合态的交换性质。}

\reading{交换振幅与占据数参照Feynman第III卷第4章\cite{feynman4}；HOM原始实验见\cite{hom}。
    下一章将粒子交换与一般的量子纠缠区别开来。}
